\Chapter{Problem description and current situation analysis}\label{chapter:problem_analysis}

\section{Project context}
\subsection{Digitalization of healthcare in the European Union}
In recent years, digitalization has become one of the main directions of healthcare system development in the European Union countries. eHealth solutions are being gradually implemented both in large hospitals and in smaller medical facilities, and electronic medical records are beginning to serve as the primary source of patient information. This process is part of the broader EU digital transformation strategy, whose goal is to increase the availability of public services and improve the quality of healthcare institutions \cite{EU_eHealth_2024}.

One of the most important assumptions of the "Digital Decade 2030" programme is to provide all European Union citizens with access to electronic medical records by 2030 \cite{EU_eHealth_2024}. Accordingly, member states are developing central systems that enable the storage and sharing of health data in digital form. According to the "2024 Digital Decade eHealth Indicator Study" report, the average level of eHealth service development in EU countries reached 79\% in 2023, which represents further growth compared with previous years \cite{EU_eHealth_2024}. Poland was identified as one of the countries with a high level of maturity in digital health services, achieving a score of 90\% \cite{EU_eHealth_2024}.

The report also shows that most European Union countries use centralized platforms for access to medical data, and an increasing number of services are available both through web portals and mobile applications \cite{EU_eHealth_2024}. This means that contemporary healthcare systems are moving away from a model based exclusively on paper documentation and are gradually transitioning to a digital environment in which information can be available in real time to authorized users.

The development of eHealth systems also stems from the need to improve the flow of information between different healthcare entities. EU reports emphasize that digital medical systems are intended to improve access to data, support continuity of care and increase patient involvement in the healthcare process \cite{EU_eHealth_2024}. Ensuring the possibility of rapid access to test results, treatment history and data concerning therapies and medications taken is becoming particularly important.

Despite the dynamic development of technology, many countries still face problems related to interoperability of systems, data integration and the availability of certain types of medical documentation \cite{EU_eHealth_2024}. The report indicates, among other things, that some data, such as medical images or information concerning implants, is still not widely available to citizens in electronic form \cite{EU_eHealth_2024}. Ensuring an appropriate level of data security, compliance with EU regulations and access exclusively for authorized persons also remains a significant challenge.

The digitalization of healthcare also involves the need to take into account the needs of various social groups. EU documents emphasize the importance of creating systems accessible to older people, people with limited digital skills and users requiring additional technical support \cite{EU_eHealth_2024}. For this reason, modern medical systems should be designed not only as technologically advanced solutions, but also as tools that are intuitive, legible and as simple as possible in everyday use.

\subsection{Project goal and assumptions}

The project concerns the creation of an integrated system supporting the management of patient information at Szpital im. Bajgla in Wrocław. The facility plans to implement a solution that enables the collection, updating and sharing of medical documentation within a single, coherent digital environment. The main goal of the project is to improve the flow of information between departments and to reduce the problems resulting from the operation of paper documentation.

The project assumption is to create a system that enables medical staff to quickly access the most important patient information, treatment history, test results and current medical orders. The solution is intended to support the daily work of doctors, nurses, laboratories and administration by organizing the flow of data and reducing the need to manually transfer documentation between departments.

The project also takes into account patients' needs related to easier access to their own medical documentation and the ability to check test results in electronic form. The system is intended to improve communication between the facility and the patient and to reduce the need to use paper documentation.

One of the key project assumptions is also to ensure a high level of security for medical data. The system should enable control of access to information, separation of user permissions and protection of data in accordance with the applicable regulations concerning privacy and information security in the European Union. An important aspect of the project also remains the creation of a solution that is intuitive and simple in everyday use, so as to reduce the problems associated with implementing new technologies in a hospital environment.
\section{Analysis of the current situation}
Szpital im. Bajgla in Wrocław is a large, multi-department medical facility. The hospital structure includes, among others, surgical departments, internal medicine, neurology, orthopaedics, oncology, paediatrics, an intensive care unit, the Hospital Emergency Department, a diagnostic laboratory, a medical imaging unit and a mortuary. Due to the scale of its activity, the facility serves a large number of patients every day, which entails the need for continuous processing and transfer of a huge amount of medical information between various hospital units.

Currently, medical documentation at Szpital im. Bajgla is kept mainly in paper form. Patient charts, test results, diagnostic descriptions, medical orders, nursing reports and administrative documents are stored physically and transferred between departments manually by medical or administrative staff. This leads to situations in which access to information is limited in time and depends on the current location of the documentation. In practice, this means that staff often have to find the necessary documents themselves, contact other departments or wait for test results and reports to be delivered.

The flow of information between departments and the laboratory and medical imaging unit is particularly problematic. Laboratory test results and descriptions of imaging examinations are transferred in a traditional manner, which extends the documentation circulation time and hinders ongoing access to current patient data. This constitutes a critical problem in situations requiring a rapid response.

The lack of a central digital system also causes problems related to the archiving of documentation. Patient data is stored in numerous binders, files and paper archives, which means that finding an older treatment history may require a large amount of time. In the case of patients hospitalized multiple times, staff must manually search documentation from previous stays in order to obtain a complete picture of the treatment course.

The current way in which the facility operates directly affects the daily work of doctors, nurses and administrative employees. A significant part of the time is devoted not to patient contact, but to organizing documentation, transferring data and completing paper forms. In many cases, the same information is rewritten several times into different documents, which increases the risk of errors and inaccuracies. The problem intensifies particularly during periods of high departmental workload or sudden situations requiring rapid access to information about a patient's health condition.

Information chaos is also visible from the patients' perspective. Hospitalized people do not have direct access to their test results or information about the course of treatment. Obtaining documentation requires contact with staff or collecting documents in paper form.

An additional challenge is the growing amount of medical data generated every day by the hospital. The development of diagnostics, a greater number of tests performed and a growing number of patients mean that the traditional model of storing documentation is becoming increasingly less efficient organizationally. The current way of working also does not correspond to the direction of development of modern healthcare in the European Union, in which increasing emphasis is placed on the digitalization of medical services, rapid access to data and integration of patient information \cite{EU_eHealth_2024}.

The analysis of the current situation shows that the main problem of Szpital im. Bajgla is the lack of a single, central environment enabling efficient management of medical documentation and information flow between users.
\subsection{Perspective of the director and the hospital organization}

The management of Szpital im. Bajgla recognizes that the current way of organizing medical documentation is no longer sufficient for a facility of such a large scale of activity. With the increase in the number of patients and the amount of tests performed, it is becoming increasingly difficult to effectively manage information flow. From the director's perspective, improving work organization and reducing situations in which medical staff devote a significant part of their time to administrative activities is particularly important.

Another problem for the hospital organization is the lack of a uniform standard for working with medical documentation. Individual departments have developed their own ways of transferring information and organizing documents, as a result of which data flow is not consistent throughout the facility.

The hospital management also draws attention to the growing requirements related to the further digitalization of healthcare in the European Union. The facility must gradually adapt to new standards concerning electronic medical documentation, data security and the availability of digital services for patients \cite{EU_eHealth_2024}.

The management expects a solution that will allow staff to obtain patient data faster, monitor the course of treatment and reduce the problems resulting from delayed documentation transfer. This is particularly important in the case of departments with high work dynamics, such as the Emergency Department or the intensive care unit, where rapid access to information can have a direct impact on patient safety.

At the same time, the management is aware that implementing a new system will be a major organizational change for the entire hospital. There is a concern that an improperly designed solution could increase the number of administrative duties or meet with resistance from some employees.
\subsection{Perspective of medical staff}
The medical staff of Szpital im. Bajgla belongs to the group of users most directly affected by the current way of organizing documentation. Doctors, nurses, paramedics, laboratory technicians and diagnostics staff participate every day in the process of transferring large amounts of information concerning patients.

From the doctors' perspective, one of the greatest problems is the lack of rapid access to the patient's complete treatment history. Obtaining all the necessary information often requires reviewing many documents, contacting other departments or waiting for test results to be delivered. Doctors point out that the time devoted to organizing documentation limits their ability to focus directly on working with the patient.

Similar problems concern nursing staff, who, in addition to duties related to patient care, are also responsible for maintaining and updating part of the medical documentation. This causes an additional organizational burden and increases the risk of mistakes resulting from illegible notes, lost documents or outdated data.

Another significant problem for staff is the flow of information concerning laboratory and diagnostic tests. Laboratory and medical imaging employees perform a large number of tests, but transferring results in traditional form extends the entire communication process. Medical staff do not have the ability to monitor the status of tests on an ongoing basis, which makes it necessary to repeatedly check the availability of results or contact other hospital departments.

Some employees are afraid of the facility's digitalization. Among the most frequently mentioned problems are concerns about the need to learn new systems, changes to existing habits and an increase in the number of administrative duties. Some employees perceive the implementation of a new solution as a potential hindrance to daily work, especially if the system were complicated. There is also a visible fear of making mistakes when using new digital tools.

At the same time, some staff recognize the benefits resulting from the introduction of a central IT system. Doctors and nurses point primarily to the need for rapid access to documentation, easier transfer of information and a reduction in the amount of paperwork. Staff expect a solution that will simplify daily duties, reduce the number of repetitive administrative activities and improve communication between departments.

From the users' point of view, the intuitiveness of the system becomes very important. The new solution should be as simple and legible as possible. The system cannot require complicated operation or a large number of additional actions when performing basic duties. The mobility of the solution is also important, particularly in the case of work in departments and when moving between patients. The possibility of rapid access to data from mobile devices could significantly improve the organization of staff work.

Employees expect a solution that will genuinely simplify the performance of duties, reduce organizational chaos and improve communication between departments, rather than another system that increases the level of complexity of work.
\subsection{Patient perspective}
Patients of Szpital im. Bajgla are the group directly experiencing the effects of problems related to the organization of documentation at the facility. The current operating model of the hospital means that patients have limited access to data concerning their own treatment and the course of hospitalization. Obtaining test results, treatment history or information about further stages of therapy most often requires contact with staff or collecting documents in paper form.
This becomes particularly problematic in the case of patients using the services of several departments simultaneously or requiring long-term treatment. Information concerning tests, consultations and medical recommendations is dispersed, as a result of which the patient often does not have a complete picture of their own treatment process. From the patients' perspective, waiting time for information is an important problem. Collecting test results or documentation often requires a personal visit to the hospital or a long wait for data to be provided by staff. In practice, this hinders the patient's contact with their own medical documentation and limits the ability to react quickly to changes related to their health condition.

Patients also point to the lack of transparency in the treatment process. Hospitalized people do not always know at what stage their tests are, when results will appear and what actions will be taken next.
At the same time, patients increasingly expect digital solutions similar to those used in other public and private services. The ability to check test results online, easier contact with the facility and faster access to medical documentation are perceived as important conveniences improving the comfort of using healthcare. Reports on eHealth development emphasize that increasing citizens' access to their own medical data is one of the main directions of modern healthcare development \cite{EU_eHealth_2024}.
\subsection{Data security and access}
One of the most important challenges related to the operation of Szpital im. Bajgla is ensuring an appropriate level of security for medical documentation and control of access to patient information. Health data belongs to the category of particularly sensitive data, which is why the manner of its storage, processing and sharing requires maintaining high security standards and compliance with applicable legal regulations.

In the current model, access to data depends on the physical location of documents. Such a model makes it difficult to effectively control who is using the documentation at a given moment and whether all information is properly protected against unauthorized access. The problem becomes particularly visible in a large facility where many people work with documentation every day.

The possibility of losing documents or improperly transferring them between departments is also a significant threat. Paper documentation may be accidentally left in the wrong place, handed to the wrong person or damaged.

From the perspective of the hospital organization, it also becomes important to separate levels of access to data. Individual employees do not need access to all information contained in medical documentation. Doctors, nurses, laboratory technicians, administration or mortuary employees use different ranges of data related to their duties. In the current paper-based model, controlling such permissions is difficult because documentation can be physically transferred between many people without the possibility of accurately monitoring access.

Requirements resulting from regulations in force in the European Union are also an important aspect. Digital healthcare systems must comply, among other things, with GDPR principles concerning personal data protection and information security \cite{EU_eHealth_2024}. This means the need to ensure appropriate user authentication, control of access to documentation and the ability to monitor actions performed in the system.
Ensuring security while maintaining simplicity of system use also remains a significant challenge. Medical staff work under significant time pressure, which is why procedures related to logging in and accessing data cannot excessively hinder daily work.
\subsection{Technical and organizational problems}
The planned digitalization of Szpital im. Bajgla involves not only the need to create a new IT system, but also a major organizational change covering the way the entire facility operates. Due to the scale of the hospital and the number of departments, implementing the new solution requires taking into account many technical, logistical and organizational problems that may affect the later operation of the system and its acceptance by users.

One of the main challenges is the need to design a system adapted to the work of many different hospital units at the same time. Each department operates somewhat differently and has its own needs related to documentation and information flow. This means that the system must be flexible enough to support various medical processes without excessively complicating the user interface.

Another technical problem is the creation of appropriate infrastructure enabling the system to function throughout the hospital. The facility requires stable network access, appropriate devices for medical staff and a secure data storage environment.
Implementing a system covering many departments means ensuring continuity of operation even in the event of failure of part of the infrastructure or technical problems.

An important challenge is also the transfer of existing paper documentation to the new digital environment. The hospital has a large amount of archival patient data that has been stored in traditional form for years. The documentation digitalization process therefore requires establishing a data migration method, setting priorities and preparing procedures that enable a gradual transition to the new working model without paralysing the facility's operations.

Preparing staff to work with the new system also remains a problem. System implementation cannot be limited solely to technical aspects. It is necessary to conduct training, prepare support materials and create support procedures for users at the initial stage of system operation.

An additional organizational challenge is the need to maintain continuity of hospital operations during the implementation of the new solution. The facility operates around the clock and cannot be temporarily shut down for system installation. In practice, this means the need for phased implementation of individual functionalities and the simultaneous operation of the old and new ways of working for a certain transition period.

Scalability of the solution also remains an important aspect. The system should enable further development of the facility, handling a growing number of patients and the possibility of expansion with new functionalities in the future.
\section{Summary of the problem analysis}
\subsection{Summary of the problem}
The analysis of the operation of Szpital im. Bajgla shows that the main problem of the facility is the lack of a central system enabling efficient management of medical documentation and information flow between departments. The current working model, based mainly on paper documentation, causes difficult access to patient data, extends the information transfer process and increases the administrative burden on medical staff. These problems affect both the organization of the hospital's work and the comfort and safety of patients.

Another significant problem is the lack of a uniform working environment for all units of the facility. The current way of organizing work also does not correspond to the growing requirements related to the digitalization of healthcare and the need for rapid access to medical data in real time.

The designed solution must take into account the needs of various user groups. The system should support the daily work of doctors, nurses, laboratories and administration, while remaining a simple and intuitive solution to use.

Ensuring an appropriate level of medical data security and creating a transparent model of access to information depending on the user's role also remain important aspects of the project.

\subsection{Target functionality}
The project will be considered successful if a system is created that enables digital handling of patient documentation at Szpital im. Bajgla without the need to use paper-based data circulation in the daily work of medical staff. The system must enable adding, viewing and updating medical documentation in real time by authorized users from various hospital departments.

A success criterion will also be ensuring central access to patient data from a single IT environment. After logging in to the system, a doctor should be able to check the patient's treatment history, laboratory test results, diagnostic documentation and current orders without the need to use additional sources of information.

The system must support various levels of access to data. Each user group should have only the permissions corresponding to the duties performed. This means, among other things, that the laboratory can add test results, a doctor has access to the patient's full medical documentation, while administration cannot view medical data not required to perform organizational duties.

Another important success criterion will be the ability for patients to access selected medical information online. After logging in, a patient should be able to check test results, visit history and basic information concerning the course of treatment without the need to collect documentation in paper form.
The project will be considered effective if the system enables simultaneous work of many users without data loss and ensures continuity of operation in a 24/7 hospital environment.
